Fix an orthogonal basis (\claimref{synthesis.completion-recovers-the-real-clifford-algebras}(iii)). For blades $e_I$, $e_J$ the geometric product is a multiple of $e_{I \triangle J}$ (\claimref{std-algebra.orthogonal-anticommute}), of grade $|I| + |J| - 2|I \cap J|$; so $e_I \wedge e_J = \langle e_I e_J \rangle_{|I|+|J|}$ is $e_I e_J = \epsilon(I, J)\, e_{I \cup J}$ when $I \cap J = \emptyset$, with $\epsilon$ the sign of the shuffle that sorts $I \cup J$, and $0$ otherwise. These are exactly the structure constants of $\Lambda V$ in the basis $e_{i_1} \wedge \dots \wedge e_{i_k}$, and they involve no value $Q(e_i)$. The linear map $\Phi$ sending that basis of $\Lambda V$ to the blade basis $e_I$ (\claimref{std-algebra.blade-basis}) is therefore bijective, grade preserving and multiplicative; so $(\Cl, \wedge)$ is associative and unital, being isomorphic to $\Lambda V$. Since $\Phi(v) = v$ and $\Lambda V$ is generated by $V$, $\Phi$ is the unique algebra map with this property, and it does not depend on the orthogonal basis.

The iterated outer product $v_1 \wedge \dots \wedge v_k$ and the antisymmetrised product $[v_1, \dots, v_k]$ are both multilinear and alternating in the $v_j$, and they agree on distinct orthogonal basis vectors, where both are $e_I$ up to the same sign (\claimref{std-algebra.grade-decomposition}); so they are equal.

The symbol map: in the exterior model, $c(e_i)$ acts on $e_J$ with $i \notin J$ as $e_i \wedge e_J$, since the contraction with $B(e_i, \cdot)$ vanishes on the other orthogonal basis vectors. Hence $e_I \cdot 1 = e_{i_1} \wedge \dots \wedge e_{i_k}$, so the symbol map agrees with $\Phi^{-1}$ on a basis. Since $\Cl(V, Q)$ is unique up to a unique isomorphism (\claimref{std-algebra.existence-uniqueness}), the statement does not depend on the model of the algebra used.
